% =============================================================
%  Einleitung -- Standalone-Dokument
%  (Nur die Introduction, ohne Abstract/Methodik/Results/Disc.)
%
%  Konkrete Ergebnis-Zahlen sind durch X / Y ersetzt, damit nicht
%  erkennbar ist, wie weit das Hauptpaper bereits ist.
% =============================================================

\documentclass[conference]{IEEEtran}
\IEEEoverridecommandlockouts

\usepackage{cite}
\usepackage{amsmath,amssymb,amsfonts}
\usepackage{graphicx}
\usepackage{hyperref}
\usepackage{url}
\usepackage{siunitx}
\usepackage[T1]{fontenc}
\usepackage[utf8]{inputenc}

\hypersetup{
  colorlinks=true,
  linkcolor=blue!50!black,
  citecolor=blue!50!black,
  urlcolor=blue!50!black,
}

\begin{document}

\title{App-Level TinyML for Smartphone Battery-Life Prediction}

\author{
  \IEEEauthorblockN{Keno Schürger}
  \IEEEauthorblockA{
    Technische Hochschule Würzburg-Schweinfurt (THWS)\\
    Email: \texttt{keno.schuerger@study.thws.de}
  }
}

\maketitle

% =====================================================================
\section{Introduction}
% =====================================================================
Predicting the remaining battery life of a smartphone is a long-standing
problem with two qualitatively different solution classes. \emph{Analytical}
methods (e.g.\ the linear extrapolation used by most legacy ``Settings
$\rightarrow$ Battery'' screens) project the recently observed drain rate
into the future. \emph{Machine-learning} methods, in contrast, attempt to
exploit context: which apps are active, how bright is the screen, what is
the battery temperature.

Since Android~12 (October~2021), the platform exposes a system
estimator via
\texttt{PowerManager.\allowbreak getBatteryDischargePrediction()}\cite{android2021powermanager},
which returns a duration estimate computed inside the operating system.
As a system process the estimator runs at a privilege level not
available to third-party apps and can in principle access hardware
counters that the Android permission model does not expose to user-%
installed applications. The companion method
\texttt{isBatteryDischargePredictionPersonalized()} indicates whether
the returned estimate has been adapted to the specific device's usage
history.

We ask: \emph{how well can a TinyML model running inside an
Android app predict remaining battery life, compared to analytical
baselines and Google's native API --- measured on accuracy and
on-device efficiency?} TinyML here refers to small quantised neural
networks (typically tens of kilobytes) designed for low-power edge
devices~\cite{banbury2021mlperf}. While the TinyML literature largely
targets microcontrollers, smartphones share two relevant properties:
deployment constraints (model size matters for APK weight, latency
matters for battery cost) and ubiquity.

This paper provides an empirical answer for the smartphone case. We
report a six-way head-to-head comparison on identical real-world data
($X$ measurements collected over an $X$-day window from four Android
devices spanning two manufacturers), explicitly controlled for a class
of data-leakage pitfalls that we found to inflate apparent results by a
data-diversity-dependent factor.

\subsection*{Contributions}
\begin{enumerate}
\item \textbf{Methodological observation:} On sliding-window time-series
data, a random shuffle of \emph{sequences} (the default
\texttt{train\_test\_split}) silently leaks future information into
training. On our data the inflation factor depends on data
diversity ($\sim$$X\times$ on the multi-device dataset,
$\sim$$Y\times$ on the single-device subset).
We propose and measure a segment-level split as the principled
alternative.

\item \textbf{Multi-device six-way comparison with statistical rigour:}
We collect data on four Android devices (Xiaomi 2107113SG, Pixel~7~Pro,
8~Pro and 9~Pro~XL) and compare TinyML against four explicit baselines
(Random Forest, mean predictor, linear drain rate, exponential fit) and
the native Google~API. Differences are reported with bootstrap 95\%
confidence intervals; pairwise significance uses a paired permutation
test for MAE and a paired bootstrap on $\Delta C$ for the $C$-index
(per-sample swap is invalid for a pairwise metric), with
Benjamini--Hochberg correction over 15 tests per family.

\item \textbf{Device-dependent learnability:} Per-device evaluation
reveals a strong hardware effect. On the Pixel~7~Pro the TinyML model
achieves $C{=}X$, several times the gap to the mean predictor that the
same model attains on the Xiaomi ($C{=}Y$). This demonstrates that the
signal available to an app-level model is dominated by the underlying
sensor stack, not by the model architecture or training budget.

\item \textbf{Model-independent feature diagnostics:} Beyond
permutation-importance on the Random Forest, we report Pearson, Spearman
and mutual-information statistics for all ten features. This separates a
model-specific failure (``Conv1D is the wrong architecture'') from a
data-level limitation (``the publicly available features only carry
enough signal on certain hardware'').

\item \textbf{Reproducible open pipeline.} All code, configuration, raw
data, intermediate artefacts, and the auto-generated report are released
in a structured repository. A single command
(\texttt{python run\_pipeline.py}) reproduces every number in this paper.
\end{enumerate}

\bibliographystyle{IEEEtran}
\bibliography{einleitung_refs}

\end{document}
